\documentclass[11pt]{article}
\usepackage[a4paper,margin=18mm]{geometry}
\usepackage[no-math]{fontspec}
\setmainfont{Latin Modern Roman}
\usepackage{amsmath,amssymb,amsthm,xfrac,xspace,hyperref,graphicx,comment}
\excludecomment{DefTralics}
\providecommand{\C}{}
\providecommand{\bibliofont}{\small}
\newcommand{\ie}{i.e.,\xspace}
\newcommand{\eg}{e.g.,\xspace}
\newcommand{\etc}{etc.\xspace}
\newcommand{\cf}{cf.\xspace}
\newcommand{\resp}{resp.\xspace}
\newcommand{\smaller}{\small}
\newcommand{\Subsection}{\subsection}
\newcommand{\Subsubsection}{\subsubsection}
\newcommand{\skpt}{\smallskip}
\emergencystretch=3em
\input{author-preamble.tex}
\begin{document}
\begin{center}{\Large Stable item 1933}\end{center}
\paragraph{Source and attribution.}
Cyril Demarche, David Harari, \emph{Duality for complexes of tori over a global field of positive characteristic}, Journal de l’École polytechnique — Mathématiques 7 (2020), publisher version, DOI 10.5802/jep.129. \href{https://jep.centre-mersenne.org/articles/10.5802/jep.129/}{Publisher article}. Authors' text: \href{https://creativecommons.org/licenses/by/4.0/}{CC BY 4.0}.
\paragraph{Editorial problem boundary.}
Prove only Theorem 4.9(a) below, for the original complex \(\mathcal C=[\mathcal T_1\to\mathcal T_2]\) of tori in degrees \(-1,0\) over a nonempty open \(U\) of the printed finite-field curve, including \(U=X\). The selected dualities pair \(H_c^3(U,\widehat{\mathcal C})\) with the completion of \(H^{-1}(U,\mathcal C)\), and \(H_c^2(U,\widehat{\mathcal C})\) with the completion of \(H^0(U,\mathcal C)\). All primes, including the characteristic, are retained. Compact support, finite-level duality, torsion/completion and five-lemma core remain in full. Other degrees and Poitou--Tate statements are context only.
\paragraph{Difficulty (editorial): H3.}
Finite-flat Artin-Mazur-Milne duality controls derived torsion realizations even at the characteristic prime. Topological compact-support dévissage, divisible-kernel limit analysis, and finite-generation Tate-module vanishing produce integral perfect dualities rather than prime-to-characteristic approximations.
\paragraph{Editorial transcription note.}
Original author language, mathematics and ordering are preserved. Publisher front matter is replaced by this independent article wrapper. Bibliography is recovered from matching cached publisher HTML, including its TeX formulas. No new proof or mathematical repair is supplied. Surrounding original results are retained for conservative context and are not extra selected corpus claims.
\clearpage
\input{author-body.tex}
\begin{thebibliography}{99}
\bibitem{BD} M. Borovoi \& C. Demarche - “Manin obstruction to strong approximation for homogeneous spaces” , Comment. Math. Helv. 88 (2013) no. 1, p. 1-54
\bibitem{Bor} M. Borovoi - “The Brauer-Manin obstructions for homogeneous spaces with connected or abelian stabilizer” , J. reine angew. Math. 473 (1996), p. 181-194
\bibitem{bkiac} N. Bourbaki - Éléments de mathématique. Algèbre commutative. Chapitre 10 , Springer-Verlag, Berlin, 2007
\bibitem{Ces} K. Česnavičius - “Topology on cohomology of local fields” , Forum Math. Sigma 3 (2015), article ID e16, 55 pages
\bibitem{cesnalms} K. Česnavičius - “$p$-Selmer growth in extensions of degree $p$” , J. London Math. Soc. (2) 95 (2017) no. 3, p. 833-852
\bibitem{demlms} C. Demarche - “Le défaut d’approximation forte dans les groupes linéaires connexes” , Proc. London Math. Soc. (3) 102 (2011) no. 3, p. 563-597
\bibitem{demphd} C. Demarche - “Suites de Poitou-Tate pour les complexes de tores à deux termes” , Internat. Math. Res. Notices (2011) no. 1, p. 135-174, short version of “Théorèmes de dualité pour les complexes de tores”, available at https://webusers.imj-prg.fr/\textasciitilde{}cyril.demarche/
\bibitem{DemH} C. Demarche \& D. Harari - “Artin-Mazur-Milne duality for fppf cohomology” , Algebra Number Theory 13 (2019) no. 10, p. 2323-2357
\bibitem{friedsus} E. M. Friedlander \& A. Suslin - “The spectral sequence relating algebraic $K$-theory to motivic cohomology” , Ann. Sci. École Norm. Sup. (4) 35 (2002) no. 6, p. 773-875
\bibitem{gonz1} C. D. González-Avilés - “Arithmetic duality theorems for 1-motives over function fields” , J. reine angew. Math. 632 (2009), p. 203-231
\bibitem{gonz2} C. D. González-Avilés \& K.-S. Tan - “The generalized Cassels-Tate dual exact sequence for 1-motives” , Math. Res. Lett. 16 (2009) no. 5, p. 827-839
\bibitem{God} R. Godement - Topologie algébrique et théorie des faisceaux , Hermann, Paris, 1973
\bibitem{hss} D. Harari, C. Scheiderer \& T. Szamuely - “Weak approximation for tori over $p$-adic function fields” , Internat. Math. Res. Notices (2015) no. 10, p. 2751-2783
\bibitem{lang} S. Lang - “Algebraic groups over finite fields” , Amer. J. Math. 78 (1956), p. 555-563
\bibitem{MilEC} J. S. Milne - Étale cohomology , Princeton Math. Series , vol. 33 , Princeton University Press, Princeton, N.J., 1980
\bibitem{MilADT} J. S. Milne - Arithmetic duality theorems , BookSurge, LLC, Charleston, SC, 2006
\bibitem{nsw} J. Neukirch, A. Schmidt \& K. Wingberg - Cohomology of number fields , Grundlehren Math. Wiss. , vol. 323 , Springer-Verlag, Berlin, 2008
\bibitem{Ros} Z. Rosengarten - “Tamagawa numbers and other invariants of pseudo-reductive groups over global function fields” , 2018, to appear in Algebra Number Theory
\bibitem{San} J.-J. Sansuc - “Groupe de Brauer et arithmétique des groupes algébriques linéaires sur un corps de nombres” , J. reine angew. Math. 327 (1981), p. 12-80
\bibitem{col} J.-P. Serre - Corps locaux , Publications de l’Université de Nancago , vol. VIII , Hermann, Paris, 1968
\bibitem{SGA4} M. Artin, A. Grothendieck \& J.-L. Verdier - Théorie des topos et cohomologie étale des schémas. Tome 3 , Lect. Notes in Math. , vol. 269, 270 \& 305 , Springer-Verlag, Berlin-New York, 1973, Séminaire de Géométrie Algébrique du Bois-Marie 1963–1964 (SGA 4)
\bibitem{sga3} M. Demazure \& A. Grothendieck - Schémas en groupes (SGA 3). Tome I. Propriétés générales des schémas en groupes , Documents mathématiques , vol. 7 , Société Mathématique de France, Paris, 2011, article ID e16, Revised and annotated edition of the 1970 original
\bibitem{SP} Stacks Project Authors - “The Stacks Project” , 2019, https://stacks.math.columbia.edu
\end{thebibliography}
\end{document}
